% ===========================================================================
\section{Information can only be lost}\label{sec:information-properties}
\warmth{0}\quad\whisper{Condition, expand, and discard a non-negative term.}

\begin{strand}
Mutual information behaves like shared evidence. It is non-negative and
symmetric, it can be revealed one variable at a time, and processing an
observation cannot manufacture information about its source.
\end{strand}

\block{The core identities}
\begin{proposition}[information properties]\label{prop:mi-properties}
For discrete random variables $X$ and $Y$,
\[
  \MI(X;Y)\ge0.
\]
Equality holds if and only if $X$ and $Y$ are independent. Moreover,
\begin{align*}
  \MI(X;Y)&=\MI(Y;X),\\
  \MI(X;Y)&=\Ent(X)-\Ent(X\mid Y)
           =\Ent(Y)-\Ent(Y\mid X).
\end{align*}
\end{proposition}

\begin{proof}
By Definition~\ref{def:mi},
\[
  \MI(X;Y)=\KL(p_{XY}\Vert p_Xp_Y).
\]
The information inequality gives
non-negativity, with equality exactly when $p_{XY}=p_Xp_Y$, which is
independence. Symmetry follows because the defining log-ratio is unchanged
when $X$ and $Y$ are exchanged. The entropy identities were obtained from the
chain rule in Proposition~\ref{prop:chain}.
\end{proof}

\begin{definition}[conditional mutual information]\label{def:conditional-mi}
The information shared by $X$ and $Y$ after $Z$ is known is
\[
  \MI(X;Y\mid Z)
  :=\sum_zp_Z(z)\KL\!\left(
    p_{XY\mid Z}(\,\cdot,\cdot\mid z)
    \,\middle\Vert\,
    p_{X\mid Z}(\,\cdot\mid z)p_{Y\mid Z}(\,\cdot\mid z)\right).
\]
Hence $\MI(X;Y\mid Z)\ge0$.
\end{definition}

\begin{proposition}[chain rule for mutual information]\label{prop:mi-chain}
\[
  \MI(X_1,\ldots,X_n;Y)
  =\sum_{i=1}^n
    \MI(X_i;Y\mid X_1,\ldots,X_{i-1}).
\]
\end{proposition}

\begin{proof}
For each term, use the conditional entropy identity
\[
  \MI(U;Y\mid V)=\Ent(Y\mid V)-\Ent(Y\mid U,V),
\]
with $U=X_i$ and $V=(X_1,\ldots,X_{i-1})$. The adjacent conditional
entropies telescope:
\begin{align*}
  \sum_{i=1}^n\MI(X_i;Y\mid X_1,\ldots,X_{i-1})
  &=\sum_{i=1}^n\bigl[
    \Ent(Y\mid X_1,\ldots,X_{i-1})
    -\Ent(Y\mid X_1,\ldots,X_i)\bigr]\\
  &=\Ent(Y)-\Ent(Y\mid X_1,\ldots,X_n).
\end{align*}
The final expression is $\MI(X_1,\ldots,X_n;Y)$.
\end{proof}

\begin{yourturn}
Expand the two-variable case:
\[
  \MI(X_1,X_2;Y)
  =\MI(X_1;Y)+
   \answerin[3.9cm]{\MI(X_2;Y\mid X_1)}.
\]
\recall{The missing term is $\MI(X_2;Y\mid X_1)$.}
\end{yourturn}

\block{Data processing}

\begin{definition}[Markov chain]
We write $X\to Y\to Z$ when $X$ and $Z$ are conditionally independent given
$Y$, equivalently
\[
  p_{Z\mid XY}(z\mid x,y)=p_{Z\mid Y}(z\mid y).
\]
\end{definition}

\begin{theorem}[data processing inequality]\label{thm:dpi}
If $X\to Y\to Z$ is a Markov chain, then
\[
  \MI(X;Z)\le\MI(X;Y).
\]
\end{theorem}

\begin{proof}
Expand the same mutual information in two orders:
\[
  \MI(X;Y,Z)
  =\MI(X;Y)+\MI(X;Z\mid Y)
  =\MI(X;Z)+\MI(X;Y\mid Z).
\]
The Markov condition gives $\MI(X;Z\mid Y)=0$, whereas conditional mutual
information is non-negative. Therefore
\[
  \MI(X;Y)
  =\MI(X;Z)+\MI(X;Y\mid Z)
  \ge\MI(X;Z).
\]
\end{proof}

\begin{corollary}[deterministic processing]
For any function $f$,
\[
  \MI(X;f(Y))\le\MI(X;Y).
\]
\end{corollary}

\begin{proof}
Because $f(Y)$ is determined by $Y$, the variables form the Markov chain
$X\to Y\to f(Y)$. Apply Theorem~\ref{thm:dpi}.
\end{proof}

\begin{yourturn}
Name the two proof moves behind data processing:
\workspace[1]
\answerprose{(1) Expand $\MI(X;Y,Z)$ in two orders. (2) Use the Markov
condition $\MI(X;Z\mid Y)=0$ and non-negativity
$\MI(X;Y\mid Z)\ge0$.}
\recall{Two moves: (1) expand $\MI(X;Y,Z)$ in two orders; (2) use
$\MI(X;Z\mid Y)=0$ and $\MI(X;Y\mid Z)\ge0$.}
\end{yourturn}
